\chapter{CONTACT}\label{ch:contact}
\markboth{Contact}{Contact}

\chapquote{``You know, there are four hundred billion stars out there, just in our galaxy alone. If only one out of a million of those had planets, and just one out of a million of those had life, and just one out of a million of those had intelligent life; there would be literally millions of civilisations out there.''}{Ellie Arroway}


\begin{center}
\vspace{-10pt}
{\itshape\color{steel} This chapter is dedicated to the life, legacy, and scientific work of the
late Professor Stephen Hawking (1942--2018).}\par
\vspace{6pt}
\adjustimage{max size={2.35in}{2.9in}}{hawking.jpg}\par\vspace{4pt}
{\RaggedRight\scriptsize\color{steel}Professor Stephen Hawking lecturing for NASA's 50th anniversary,
21 April 2008. Photograph: NASA/Paul Alers, public domain.\par
\vspace{1pt}{\color{midgrey} Image sourced from:} \srclink{https://images.nasa.gov/details/200804210005HQ}\par}
\end{center}

\clearpage

This chapter belongs to three men before it belongs to any argument, so they go first.

Frank Drake and Seth Shostak between them account for most of what can honestly be called
\emph{science} in the question of contact. One of them invented the field's method and its
arithmetic; the other has spent three decades defending both in public, usually against people
who would rather hear something more exciting. Neither is a Ufologist. Neither has ever claimed a
detection. That is exactly why they are the right voices to open with: in a subject where the
standard career move is to announce something, these two built the only programme in the territory
that says in advance what would count as a discovery and then reports, year after year, that it has
not happened.

\medskip\noindent\textbf{Frank Drake.}
Drake (28~May 1930 --- 2~September 2022) was born in Chicago and, by his own account, worked out
as a child of eight that other stars might have inhabited planets --- and then kept quiet about it,
because it was not the sort of thing one said. He read engineering physics at Cornell on a Navy
scholarship, heard Otto Struve lecture on the likelihood that planets are common, served his time as
an electronics officer at sea, and took his doctorate in radio astronomy at Harvard in 1958. Then he
went to the new National Radio Astronomy Observatory at Green Bank, West Virginia, and did the thing
nobody else had done: he pointed the instrument.

Project Ozma, in the spring and early summer of 1960, was the first modern search for
extraterrestrial intelligence --- roughly 150 hours on the 85-foot Tatel dish, aimed at two nearby
Sun-like stars, Tau Ceti and Epsilon Eridani, listening near the 1420\,megahertz hydrogen line
(Section~\ref{sec:seti}). It found nothing, apart from one heart-stopping false alarm that turned out
to be terrestrial. Drake published that null result, and in doing so established the discipline's
single most valuable habit.

The following year he had to write an agenda. In November 1961 a dozen people met at Green Bank to
discuss whether the search was worth pursuing at all --- Struve in the chair, Philip Morrison,
Bernard Oliver, Melvin Calvin (who learned mid-meeting that he had won the Nobel Prize), a
twenty-seven-year-old Carl Sagan, and the rest of the group that afterwards called itself the Order
of the Dolphin. Drake listed on a blackboard every quantity you would need in order to estimate how
many civilisations were currently detectable, and noticed that the quantities multiplied. That
blackboard is the Drake Equation, and I take it apart properly ---\,including what is wrong with
it\,--- in Section~\ref{sec:drake}.

What he did with the rest of his life matters for this chapter specifically. He went to Cornell,
ran the Arecibo Observatory and later the National Astronomy and Ionosphere Centre, designed the
Pioneer plaques with Sagan in 1972, drafted the Arecibo message of 1974 (Section~\ref{sec:mathlang}),
contributed to the Voyager records of 1977 (Section~\ref{sec:messagesfuture}), moved to the
University of California, Santa Cruz as dean of natural sciences, and helped found the SETI Institute
in 1984, chairing its board for years afterwards. He died at home in Aptos, California at
ninety-two. \hi{Every single artefact and argument in this chapter is downstream of him.} He is also,
not incidentally, the man who both listened \emph{and} transmitted, which makes his position in
Section~\ref{sec:firstcontactdebate} more interesting than either camp usually admits.

\medskip\noindent\textbf{Seth Shostak.}
Shostak (born 20~July 1943, Arlington, Virginia) says he was hooked at ten by a library book full
of photographs of galaxies. He read physics at Princeton, took a doctorate in astronomy at Caltech in
1972 measuring neutral hydrogen in galaxies, spent a decade doing radio astronomy at Groningen in the
Netherlands, and then --- a detail I like, because it is the opposite of the usual career of a
professional believer --- left research for several years to run a computer graphics company in
California. He joined the SETI Institute in the early 1990s and has been its senior astronomer, and
far more importantly its public voice, ever since: thousands of lectures, a weekly science radio
programme, \emph{Sharing the Universe} in 1998, \emph{Confessions of an Alien Hunter} in 2009,
and testimony before the US House Committee on Science, Space and Technology in May 2014.

His appearances in this book so far have been adversarial. He is the man in
Section~\ref{sec:ufologist} who has spent thirty years patiently explaining why UFO reports are not
evidence of visitors, and I noted there that his negative case rests on an argument from
expectation. All of that stands. But it is worth being clear about \emph{why} he does it, because
the common accusation --- that he is a professional debunker with an institutional interest in
denying aliens --- gets it exactly backwards. Shostak runs a search programme for extraterrestrial
intelligence. He believes we will find one. He told Congress he expected evidence within two dozen
years, and he was not being coy about it: he has a bet, a date, and a falsifiable claim, which is
more than anybody else in this field has ever put on the table.

His reason is arithmetic. SETI's search speed --- how many star systems can be examined, across how
many frequency channels, per unit of telescope time --- has been doubling every eighteen months or so
for decades, because it rides digital signal processing rather than steel. Extrapolate that curve and
the first million-star survey arrives in the 2030s. \hi{That is not a prophecy, it is a projection, and
it will be wrong or right on a schedule.} You will notice that this is the same manoeuvre Drake made
in 1961: take an unanswerable question, decompose it into quantities, and force the argument to happen
in numbers where it can be checked.

\medskip\noindent\textbf{Hawking's arithmetic.}
The dedication at the head of this chapter is not decorative and it is not a change of subject.

Stephen Hawking spent his career on questions about the universe that could not be answered by going
and looking --- the interiors of black holes, the first instant of the Big Bang, the boundary
conditions of spacetime --- and he answered them, or advanced them, using mathematics and rigour
alone. That is the single most relevant skill set in this book. He also said two things about
extraterrestrial life that are usually reported as a contradiction and are in fact a single coherent
position. He said the numbers make alien life plausible. And he said that if it exists, we should
probably not advertise ourselves, because ``if aliens ever visit us, I think the outcome would be
much as when Christopher Columbus first landed in America, which didn't turn out very well for the
Native Americans.''

That is not pessimism. It is the application of the only data point we have --- what happens on this
planet when a technologically advanced culture encounters a less advanced one --- to a question where
we have no other data at all. Hawking was doing what I have been asking you to do for seventeen
chapters: reasoning carefully from what is actually known, refusing to fill gaps with what he
preferred to be true, and being willing to say something unwelcome. He was also, and this matters for
a man who lost the use of his voice, obsessed with the problem of communication. A chapter about
whether we could ever talk to something not from here belongs to him.

And the two positions dovetail more precisely than the obituaries suggested. Hawking's first claim ---
``to my mathematical brain, the numbers alone make thinking about aliens perfectly rational'' --- is
Drake's equation restated in a sentence. It is the argument that with a hundred billion galaxies of a
hundred billion stars apiece, and planets now known to be ordinary rather than rare, a prior
probability of zero is the least defensible number available. That is precisely the case Drake wrote on
a blackboard and Shostak has spent thirty years turning into telescope time. Hawking's second claim
--- that we should therefore listen rather than shout --- is not a retraction of the first. It is a
separate calculation about \emph{asymmetric} risk, and it lands squarely on SETI's side of the METI
argument in Section~\ref{sec:firstcontactdebate}: search hard, publish everything, reply only by
agreement.

\hi{The unifying idea across all three men is that a question with no data is not the same as a question
with no method.} You decompose it. You write down what you would need to know, you say in advance what
would count as an answer, you build the instrument that could produce one, and you accept the number
you get. Hawking did that to black holes, which nobody can visit. Drake did it to the galaxy's
population, which nobody can census. Shostak does it to the search itself, which is why his estimates
come with dates attached. Everything else in this chapter --- the barriers, the scales, the debates,
the golden record --- is an attempt to hold that standard in territory that does not naturally reward
it.

% ---------------------------------------------------------------
\section{Barriers of Communication}\label{sec:barrierscommunication}

Chapter~\ref{ch:language}'s opening argument was that we cannot reliably communicate about this
subject \emph{with each other}. Section~\ref{sec:fallacylanguage} was about the fallacy of language: that ``UFO'',
``alien'', ``sighting'' and ``proof'' mean different things to different people in the same
conversation, and that a large fraction of the disagreement in this field is two people using one
word for two concepts. If that is true among humans who share a species, a planet, a sensory
apparatus and a hundred thousand years of common cognitive history, then the standard optimistic
assumption of contact literature --- that a sufficiently intelligent species will of course be able
to make itself understood --- deserves to be taken apart.

So let us stack up the barriers properly, from the tractable to the possibly fatal.

\textbf{The physical barrier.} Interstellar distance is the first and most underrated. The nearest
star system is 4.24 light years away, which makes the fastest possible exchange --- question, answer
--- an eight-and-a-half-year conversation, and that is the best case in the entire galaxy. A dialogue
with something 500 light years out is not a dialogue; it is two monologues separated by a
millennium, conducted between civilisations neither of which can be certain the other still exists.
Every scenario in this chapter has to be read with that constraint in place. Section~\ref{sec:rio}
is partly about the consequences.

\textbf{The sensory barrier.} We assume vision and hearing because we have them. There is no reason
to. A species that navigates by magnetic field, electrical sense, chemical gradient or echolocation
would construct a description of the world with different primitives in it. Consider that on this
planet, right now, the mantis shrimp has somewhere between twelve and sixteen photoreceptor types
against our three, and can see circular polarised light --- a property of light that no human being
has ever perceived and never will. We share four billion years of evolution with that animal and we
cannot fully imagine its visual field. Now remove the shared evolution.

\textbf{The cognitive barrier.} Human languages are all built on a small set of shared
foundations: discrete symbols, recursive grammar, roughly linear time, agent-and-object structure,
and a theory of other minds. \hi{These are not logical necessities; they are the design features of one
particular brain.} A distributed or colonial intelligence might have no concept of a singular first
person. Something with no mortality might have no tense system worth the name. And any species
without a theory of mind --- without the ability to model what another entity knows and does not know
--- cannot construct a message aimed at a stranger at all, because the entire act of explaining
requires imagining ignorance.

\textbf{The referential barrier.} This is the deep one, and it is the reason linguists are usually
more pessimistic about contact than physicists are. Translation between human languages works because
we share a world: I can point at water. Quine's famous problem --- he called it the indeterminacy of
radical translation --- is that if a stranger says ``gavagai'' when a rabbit runs past, you cannot
determine from any amount of observation whether it means \emph{rabbit}, \emph{dinner},
\emph{undetached rabbit part}, \emph{fleeting rabbit-stage}, or \emph{look out}. On Earth you resolve
this by living alongside the speakers for a year. Across 500 light years, with no shared referents
and no possibility of pointing, there is no obvious resolution procedure at all.

\begin{funfact}
We have a live control experiment for this and it is not going well. The Linear A script of Minoan
Crete, the Indus Valley script, and Rongorongo from Easter Island remain undeciphered. These were
written by \emph{human beings}, on this planet, with human bodies, human needs and human concerns,
and in the Indus case across a city civilisation of over a million people --- and we cannot read a
word. The failure is not one of intelligence or effort. It is the absence of a bilingual key. The
Rosetta Stone is the reason we can read Egyptian; nobody has ever found the interstellar equivalent,
and there is no reason to expect one to exist.
\end{funfact}

It is worth noting that Shostak accepts most of this list and thinks it changes nothing, because his
objective is not conversation. His formulation is that detection and comprehension are different
projects, and only the first one is currently on the table: a narrowband carrier at 1420\,megahertz
tells you somebody is there, whatever else it fails to tell you, and that single bit of information is
the most consequential thing science could presently deliver. \hi{The barriers in this section are
objections to a dialogue, not to a detection.} I think that is right, and I also think it is the reason
SETI and the rest of this field talk past each other so reliably --- one of them is trying to establish
a fact and the other is trying to have a relationship.

\textbf{The intentional barrier.} Finally, and least discussed: successful communication requires
both parties to want it. A species could be aware of us and uninterested, in the way that you are
aware of the ant colony in the garden and have never attempted to negotiate with it. Every barrier
above assumes good faith and mutual effort. \hi{That assumption is doing an enormous amount of unnoticed
work in almost everything ever written about first contact.}

% ---------------------------------------------------------------
\section{The Universal Language of Mathematics}\label{sec:mathlang}

The standard answer to everything in the previous section is mathematics, and it is a good answer.
It is just not as good as it is usually claimed to be.

The argument runs: physical law is the same everywhere, therefore any technological species must have
discovered the same physics, therefore we share a body of content even if we share nothing else. Any
culture capable of building a radio telescope knows that hydrogen is the most abundant element, knows
the prime numbers, knows the value of $\pi$, and knows the 21-centimetre hyperfine transition line of
neutral hydrogen at 1420 megahertz. That last one is why SETI listens where it listens: it is a
frequency any radio astronomer anywhere would already be pointing an instrument at.

\figwrap[r]{0.38}{The Arecibo message arranged as 73 rows of 23 columns --- the only sensible factorisation of 1,679.}{https://commons.wikimedia.org/wiki/File:Arecibo_message.svg}{fig:arecibo}
This is the reasoning behind every message humanity has actually sent. The \textbf{Arecibo message}
of 1974, drafted principally by Frank Drake with Carl Sagan, was 1,\!679 bits transmitted towards the globular cluster M13 (see Figure~\ref{fig:arecibo}). The number is the key: 1,679 is the product of two primes, 23 and 73, so it can
be arranged into a rectangular grid in only one sensible way, and doing so produces a bitmap showing
the numbers one to ten in binary, the atomic numbers of the elements in DNA, the structure of the
double helix, a human figure with its height, the population of Earth, a diagram of the solar system,
and a sketch of the transmitting dish. The \textbf{Pioneer plaques} of 1972 and 1973 used a different
key --- the hyperfine transition of hydrogen as a universal ruler and clock, and a map of the Sun's
position relative to fourteen pulsars, each pulsar's period written in binary. That map is
simultaneously a location and a timestamp, since pulsar periods decay predictably, so it will tell a
finder both where and \emph{when} it was made.

\pullquote[Martin Rees]{``If we ever establish contact with intelligent aliens living on a planet
around a distant star \dots\ they would be made of similar atoms to us. They could trace their
origins back to the big bang 13.7 billion years ago, and they would share with us the universe's
future. However, the surest common culture would be mathematics.''}

Those are elegant pieces of engineering. Now the objections.

First, mathematics gives you a shared vocabulary of nouns, not a language. \hi{Primes and hydrogen let
you establish that you are both technological.} They do not let you express \emph{we came in peace},
\emph{do not do that}, or \emph{we would like to negotiate}. There is no prime number for regret.
Everything humans actually need to say to a stranger --- intention, permission, warning, obligation
--- lives in exactly the part of language that physics does not underwrite.

Second, the notation is not universal even where the content is. Base ten is a consequence of having
ten fingers. Left-to-right reading, our convention that a bitmap's origin sits top-left, the very
idea of encoding a picture as a raster of pixels --- all of these are cultural. When Drake circulated
a test version of the Arecibo-style message to distinguished scientists before transmission, including
some Nobel laureates, essentially none of them decoded it correctly. Human experts, told in advance
that a message was present and that it was decodable, mostly failed. That result is treated as an
amusing anecdote and I think it is the most important datum in the entire discipline.

Third, our own recent history has quietly complicated the picture. When we built machines to find
structure in language without being told the rules, they found it --- but the representations they use
internally are not human-readable, and we cannot fully explain what they encode. We have manufactured
systems that manipulate our own language competently and whose internal states we cannot translate.
If mutual intelligibility were as automatic as the optimistic case suggests, that should not be
possible.

\pullquote[Attributed to Galileo Galilei]{``Mathematics is the language in which God has written the
universe.''}

Drake, to his considerable credit, never oversold this. He described the Arecibo transmission as a
demonstration of the new Arecibo instrument rather than a serious attempt at correspondence, pointed
out that M13 is 25,000 light years away and will not be in that part of the sky when the signal
arrives, and said plainly that the message was a proof of capability. Shostak makes the sharper version
of the same point: mathematics is the \emph{envelope}, not the letter. What SETI is actually listening
for is not content at all but artificiality --- a spectral line too narrow for nature to have made ---
and the decoding problem, which is the whole of this section, would begin only after the answer to the
first question had already changed history.

Where does that leave the idea? I think mathematics is necessary and insufficient: the only plausible
foundation for a first exchange, and nowhere near enough for a conversation. It gets you to ``we are
both here and we both count''. Everything past that is unsolved.

% ---------------------------------------------------------------
\section{The Rio Scale}\label{sec:rio}

Suppose tomorrow somebody detects something. How excited should the rest of us be?

This is a real operational problem, not a rhetorical one, and the field's answer to it is the best
piece of intellectual hygiene anywhere in this subject. The \textbf{Rio Scale} was proposed in 2000
by Iv\'an Almár and Jill Tarter, deliberately modelled on the Torino scale for asteroid impact risk
and the Richter scale for earthquakes, to give any claimed detection of extraterrestrial intelligence
a single number from 0 to 10 indicating its actual significance. A revision, Rio 2.0, was published
by Duncan Forgan and colleagues in 2018. Tarter, it is worth remembering, was the person Drake and
Shostak both handed the institute to --- she ran the observational programmes at the SETI Institute for
decades --- so the scale comes from inside the tradition described at the start of this chapter, and it
exists because the people who most wanted a detection were the people most worried about announcing a
false one.

The score is built from a few independent judgements multiplied together. How consequential would the
signal be if genuine --- a beacon aimed at us, an omnidirectional leakage signal, an artefact in the
solar system, an unambiguous physical trace? How far away is the source --- within the solar system,
nearby stars, across the galaxy, extragalactic? And, crucially, how credible is the detection ---
has it been independently confirmed by another instrument, is it a single unrepeated event, is the
instrumentation understood, has interference been excluded?

That last term is what makes the scale useful, because it forces the two questions that this book
keeps insisting on --- \emph{how important would this be} and \emph{how sure are we} --- to be scored
\emph{separately} and then combined. Almost every failure of reasoning in Chapter~\ref{ch:language}
is a case of importance leaking into confidence: because a claim would matter enormously if true,
people treat it as better supported than it is.

Figure~\ref{fig:rioscale} sets out both halves of the scale: the published 0--10 bands along the
top, and underneath them the grid the two independent judgements produce when they are plotted
against each other. Read the grid rather than the number. Every candidate in the history of the
field sits in the same place on it --- high on the vertical axis, because the consequence of a
genuine detection is enormous, and pressed hard against the left-hand edge, because the credibility
of the detection has never survived follow-up. Nothing has ever failed the consequence test. The
scale exists to stop that fact from being mistaken for progress.

\Needspace*{4.2in}
\begin{center}
\refstepcounter{figph}\label{fig:rioscale}%
\begin{figpanel}[width=0.96\textwidth]
\centering
{\setlength{\fboxsep}{0pt}\setlength{\fboxrule}{0.5pt}%
\fcolorbox{ufogreen}{white}{\adjustimage{max size={0.90\textwidth}{4.7in}}{dg_rioscale.png}}}\par\vspace{3pt}
{\RaggedRight\scriptsize\color{steel}\textbf{\textcolor{ufogreenink}{Figure \thefigph.}}
The Rio Scale. Top: the published 0--10 significance bands of Rio~2.0. Bottom: the same scale drawn
as the grid it actually is, with consequence on one axis and credibility on the other, and the five
candidates discussed in this section placed on it. The band labels and the two-term structure are
from the published scale; the position of each individual candidate is my own reading of the
evidence and not an official score.\par
\vspace{1pt}{\color{midgrey} Author's own diagram, after Alm\'ar \& Tarter (2000) and Forgan et
al., \emph{Rio 2.0} (2018).}\par}
\end{figpanel}
\end{center}
\vspace{6pt}

\figwrap[l]{0.38}{Ehman's printout of 15 August 1977, with the sequence 6EQUJ5 circled and annotated.}{https://commons.wikimedia.org/wiki/File:Wow_signal_printout.jpg}{fig:wowsignal}
Some worked examples. The \textbf{Wow! signal} of 15 August 1977 --- Jerry Ehman's 72-second
narrowband burst at the Ohio State Big Ear observatory, close to the hydrogen line, which he circled in red pen and annotated ``Wow!'' (see Figure~\ref{fig:wowsignal}) --- is the most famous candidate in history and scores low, because
despite decades of follow-up it has never recurred and never been independently confirmed. Repetition
is not a nice-to-have; it is the whole thing. \textbf{BLC-1}, a candidate found by Breakthrough Listen
in 2020 apparently coming from Proxima Centauri, would have been extraordinary given the distance,
and turned out to be terrestrial radio interference. \textbf{Oumuamua}, the interstellar object of
2017 whose unusual acceleration prompted Avi Loeb to publicly suggest an artificial origin, would
score appreciably on consequence and low on credibility, because outgassing of volatiles remains an
adequate explanation and the object is gone. And \textbf{Tabby's Star} --- KIC 8462852, whose bizarre
irregular dimming produced serious published discussion of a Dyson-sphere-like structure --- is now
best explained by an uneven dust cloud, because the dimming turned out to be wavelength-dependent and
solid objects block all colours of light equally. That last one is a genuinely beautiful piece of
reasoning and it is worth carrying with you: a single well-chosen measurement retired the exotic
hypothesis, and no argument about plausibility could have done it.

The most instructive case is the most recent one, because it happened in public and in real
time. \textbf{3I/ATLAS}, discovered on 1 July 2025 by the ATLAS survey telescope in Chile, is
only the third interstellar object ever identified, after 'Oumuamua and the comet 2I/Borisov.
It was large, it was old --- the trajectory points to the Milky Way's thick disc, which makes it
plausibly older than the solar system itself --- and it passed perihelion in late October 2025
on the far side of the Sun from Earth. Avi Loeb, the same astrophysicist who had argued the
artificial case for 'Oumuamua, published a running list of what he called anomalies, which grew
from a handful to a dozen and then to fifteen: the alignment of the orbit with the ecliptic, the
close approaches to Mars, Venus and Jupiter but not Earth, a jet pointing sunward, nickel vapour
without the accompanying iron, non-gravitational acceleration, unusual polarisation. His stated
position was carefully hedged --- that the object \emph{might} be technological and that the
possibility deserved observation rather than ridicule.

Score it properly and the exercise is short. Consequence: as high as it gets, an artefact from
another system inside our own. Credibility: low, and it fell as the data came in. Each anomaly
had a cometary explanation available, several of them published within days by other
astronomers; the geometry of the close approaches is a selection effect of an orbit that happened
to lie near the plane where the planets are; Breakthrough Listen pointed radio telescopes at it
and found nothing but a comet with a coma. The genuinely strange result, when it arrived in the
spring of 2026, was chemical rather than technological: an extraordinary ratio of heavy water,
higher than anything measured in our own system, which tells us something real about the
conditions where the object formed. By March 2026 it had made its closest approach to Jupiter
and was on its way out, fading, and will not return.

I include it because it is the cleanest available demonstration of the two things this chapter is
for. First, the method worked exactly as intended: an unusual object was flagged, the exotic
hypothesis was stated out loud, the observations were taken, and the hypothesis lost on the
evidence rather than on authority. Second, \hi{a long list of individually unexplained details is
not the same thing as one confirmed measurement, and stacking anomalies raises the consequence
score without moving the credibility score at all.} Loeb's insistence that the question be asked
is defensible and, in my view, useful. The tally of fifteen anomalies is not evidence, and the
grid in Figure~\ref{fig:rioscale} is the reason: 3I/ATLAS sits with the rest of them, high on
consequence and hard against the left-hand wall.

Shostak's public conduct through several of these episodes is the reason I take him seriously. He was
the man the news programmes rang for a quote, he had an institutional interest in excitement, and his
line was consistently the deflating one: interference until proven otherwise, dust until proven
otherwise, no announcement without a second telescope. \hi{An advocate who talks down his own field's
best headlines is demonstrating exactly the discipline the rest of this subject lacks.}

\begin{funfact}
There is also a formal etiquette for what to do afterwards. The International Academy of Astronautics
maintains a ``Declaration of Principles Concerning Activities Following the Detection of
Extraterrestrial Intelligence'' --- the SETI post-detection protocols --- which say, in essence:
verify before announcing, tell the astronomical community and the United Nations, publish the data
openly, and \emph{do not reply} on behalf of humanity without international consultation. It has no
legal force whatsoever. It is a gentlemen's agreement about the most consequential press release in
history, and the honest assessment is that in the actual event it would last about forty minutes
before somebody leaked it.
\end{funfact}

% ---------------------------------------------------------------
\section{First Contact, The Debate}\label{sec:firstcontactdebate}

On 21 September 1987, standing at the rostrum of the United Nations General Assembly in New York,
the President of the United States told the assembled delegations of the world that their
differences would evaporate overnight if something arrived from outside the solar system. He was
not speaking off the cuff, and he was not speaking to the room. He was speaking to one man in
Moscow.

\pullquote[Ronald Reagan, address to the 42nd Session of the United Nations General Assembly, New York, 21 September 1987]{In our obsession with antagonisms of the moment, we often forget how much unites all the members of humanity. Perhaps we need some outside, universal threat to make us recognise this common bond. I occasionally think how quickly our differences worldwide would vanish if we were facing an alien threat from outside this world.}

The remark is usually quoted as a curiosity. It is better read as diplomacy. Reagan had already
put the same proposition to Mikhail Gorbachev privately at the Geneva summit of November 1985, and
Gorbachev confirmed it publicly in a speech at the Grand Kremlin Palace on 16 February 1987 ---
seven months \emph{before} the United Nations address:

\pullquote[Mikhail Gorbachev, Grand Kremlin Palace, Moscow, 16 February 1987]{At our meeting in Geneva, the U.S.\ President said that if the earth faced an invasion by extraterrestrials, the United States and the Soviet Union would join forces to repel such an invasion. I shall not dispute the hypothesis, though I think it is early yet to worry about such an intrusion.}

So the two men who between them controlled roughly sixty thousand nuclear warheads used a
hypothetical first contact as a shared thought experiment about cooperation, and then one of them
said it out loud in the one chamber on Earth where every nation has a seat. That is the correct
place to begin this section, because it establishes the frame the rest of it needs. First contact
is not, institutionally speaking, an astronomy problem. It is a governance problem, and the venue
for governance problems involving all of humanity is the United Nations.

And it is worth noticing what Reagan did \emph{not} say. He did not announce a programme, a
treaty, a working group, or a plan. Neither superpower tabled anything. The most powerful
rhetorical use of first contact in the history of the General Assembly produced exactly zero
institutional output, and that gap between rhetoric and machinery is the subject of everything
that follows.

\subsection{The Question Worth Asking}

Almost every popular treatment of governments and this subject asks whether officials are secretly
in contact, secretly negotiating, or secretly reverse-engineering hardware. That question cannot be
answered from open sources, which is precisely why it is popular: an unanswerable question never
embarrasses the person asking it. There is a different question, it is answerable from the public
record, and the answer turns out to be far more damning than any conspiracy theory.

\actualq{are the governments working with aliens?}{what are governments actually DOING to prepare
us for first contact?}

That question has a measurable answer. Preparation leaves a paper trail: statutes, budget lines,
named offices, reporting channels, published protocols, staffed agencies, declassification
schedules. You can count those things. And when you count them, what you find is that the world
has built a respectable amount of machinery for \emph{detecting} anomalies and almost none for
\emph{responding} to a confirmed detection. The rest of this section is that audit, taken
institution by institution, and the reader should keep score.

\subsection{What the United Nations Has Actually Done}

The UN's formal engagement with this subject is small, real, and almost entirely the work of one
country. Grenada, under Prime Minister Sir Eric Gairy, pushed the item onto the agenda at the 32nd
and 33rd sessions of the General Assembly. On 18 December 1978, at its 87th plenary meeting and on
the recommendation of the Special Political Committee, the Assembly adopted Decision 33/426 as a
consensus text. It remains the only formal action the General Assembly has ever taken on the
subject, so it is worth being precise about what it says.

Decision 33/426 does four things. It takes note of Grenada's statements and draft resolutions. It
\emph{invites} interested Member States to coordinate scientific research and investigation into
extraterrestrial life, including unidentified flying objects, at the national level, and to inform
the Secretary-General of their observations, research and evaluations. It requests the
Secretary-General to transmit Grenada's material to the Committee on the Peaceful Uses of Outer
Space (COPUOS) for its 1979 session. And it grants Grenada the right to address that Committee.

\hi{Read the verbs. Nothing in Decision 33/426 is binding, funded, or staffed.} Grenada had asked
for an agency or department of the United Nations to undertake, coordinate and disseminate research;
what it received was an invitation to Member States to do it themselves and write in. The
difference between those two outcomes is the entire history of global governance on this issue in
microcosm. Gairy was deposed in a coup in March 1979 while abroad, the Grenadian initiative died
with his government, and COPUOS has never established a standing item on the subject since.

\begin{funfact}
\funfactlead{The 1978 briefing nobody remembers.} Gairy did not arrive empty-handed. In support of
the 1978 item he brought a delegation that included the astronaut Gordon Cooper, the astronomer
J.\,Allen Hynek, the scientist Jacques Vall\'ee, and the researcher Stanton Friedman --- arguably
the most credentialled panel ever assembled on this topic in a UN building. Contemporaneous US
diplomatic cabling from the mission in New York recorded the briefing in dry, faintly weary prose.
The panel was heard. Nothing was built.
\end{funfact}

One later episode deserves correction, because it is repeatedly misreported. In 2010 the British
press announced that Mazlan Othman, then Director of the UN Office for Outer Space Affairs
(UNOOSA), had been designated humanity's point of contact for extraterrestrial communication ---
the ``alien ambassador'' story. Othman denied it flatly and with some humour: it sounded very
cool, she said, but she had to deny it. UNOOSA has no such mandate, no such post, and no such
budget line. \hi{The single most widely believed fact about UN preparation for first contact is
false, and the truth is that the post does not exist.}

What the UN system does have is adjacent and genuinely useful. The Outer Space Treaty of 1967
obliges States Parties to avoid harmful contamination of celestial bodies and adverse changes in
Earth's environment resulting from the introduction of extraterrestrial matter --- Article IX, the
nearest thing in international law to a biosecurity clause for a returning sample or a landing.
COPUOS and its Legal Subcommittee remain the standing forum where such questions would have to be
negotiated. The Committee on Space Research (COSPAR) maintains the planetary-protection policy that
actually governs sterilisation of outbound spacecraft and containment of inbound samples. The
International Asteroid Warning Network and the Space Mission Planning Advisory Group, both endorsed
through COPUOS in 2013--2014, are working models of exactly the thing that is missing here: a
standing international body for detecting an external event, verifying it, and coordinating a
response. \hi{Humanity has built a functioning global warning system for rocks. It has not built one
for a signal.}

\subsection{The Global Ledger, Condensed}

With the UN frame established, the practical question becomes what individual states are doing,
because Decision 33/426 pushed the work down to the national level and left it there. Dozens of
countries have some history here and a complete narrative would run to a book of its own, so the
following tables condense it. Each entry is a real institution or instrument, with the years it has
operated and, in the final column, the only assessment that matters for this section: does it
prepare the public for anything?

\begin{center}\small
\textbf{\textcolor{ufogreenink}{Standing national bodies: detection and investigation}}\par\vspace{4pt}
\begin{tabular}{@{}p{0.20\textwidth}p{0.30\textwidth}p{0.40\textwidth}@{}}
\toprule
\textbf{State} & \textbf{Instrument} & \textbf{Preparation value}\\
\midrule
United States & AARO, the All-domain Anomaly Resolution Office, created by statute in 2022 with
annual reporting duties to Congress & Highest in the world for data collection and pilot reporting;
no public-communication or post-confirmation mandate\\[3pt]
France & GEIPAN at CNES, continuous since 1977 (as GEPAN, then SEPRA) & The gold standard for
transparency: case files published online with classifications, including the unexplained ones\\[3pt]
Uruguay & CRIDOVNI, inside the Air Force since 1979 & The longest continuously operating military
investigation body in South America; procedural, not public-facing\\[3pt]
Chile & CEFAA under the civil aviation authority, 1997--present & Aviation-safety framing, works
with pilots and controllers; no contact protocol\\[3pt]
Peru & DIFAA in the Air Force, 2001, reactivated 2013 & Investigative only\\[3pt]
Canada & Sky Canada Project, Office of the Chief Science Advisor, report published 2025 with
fourteen recommendations, including a federal office to receive public and pilot reports & The
clearest recent statement anywhere that the reporting system itself, not the phenomenon, is the
thing that needs fixing\\
\bottomrule
\end{tabular}
\end{center}

\begin{center}\small
\textbf{\textcolor{ufogreenink}{Legislatures, archives and stated policy}}\par\vspace{4pt}
\begin{tabular}{@{}p{0.20\textwidth}p{0.30\textwidth}p{0.40\textwidth}@{}}
\toprule
\textbf{State} & \textbf{Instrument} & \textbf{Preparation value}\\
\midrule
United Kingdom & MoD desk closed 2009; files released through the National Archives in batches from
2008 & Retrospective transparency, no forward-looking function\\[3pt]
Japan & Defence Ministry procedure for aircrew encounters, ordered by the Defence Minister in 2020,
and active Diet questioning & Rare example of a written armed-forces procedure for the encounter
itself\\[3pt]
Brazil, Argentina, Mexico, New Zealand, Sweden, Denmark, Australia & Declassified national
archives, air-force study groups, or closed historical projects & Historical record only\\[3pt]
China & FAST radio telescope conducts SETI observation programmes & Detection capability without a
published response policy\\[3pt]
Russia & No standing civil body; Soviet-era military collection programmes documented
retrospectively & None\\[3pt]
European Union & No competence claimed; individual member states act alone & None\\
\bottomrule
\end{tabular}
\end{center}

Read the right-hand columns of both tables together and the pattern is unmistakable. Almost every
entry is an instrument for \emph{collecting reports}, a handful are instruments for
\emph{releasing history}, and not one of them is an instrument for \emph{telling the public what to
do if a detection is confirmed}. \hi{Governments have built inboxes. Nobody has written the letter
that would be sent out.}

\subsection{The Protocols That Do Exist, and Their One Fatal Flaw}

It is not true that nothing has been drafted. In 1989 the International Academy of Astronautics
adopted the \emph{Declaration of Principles Concerning Activities Following the Detection of
Extraterrestrial Intelligence} --- the SETI post-detection protocols. They are sensible and short.
Verify before announcing. Tell other observers so the detection can be independently confirmed.
Inform the scientific community, the Secretary-General of the United Nations, and the public,
promptly and openly. Preserve the data. Do not reply on behalf of humanity without international
consultation. A revised and expanded successor has been worked on since, and in 2022 the University
of St Andrews established the SETI Post-Detection Hub, which now coordinates more than fifty
researchers across science, governance, communication and society precisely because the 1989
document is too thin for the job.

The flaw is structural rather than intellectual. \hi{The SETI protocols have no legal force
whatsoever.} They were adopted by a learned academy, not by a treaty body; they bind nobody, they
fund nothing, they cannot compel a government to disclose and cannot stop a private radio astronomer
from transmitting. Alexander Zaitsev sent messages from the Yevpatoria dish in Ukraine without
asking anyone, which settles the point empirically. The nearest thing humanity has to a first-contact
plan is a voluntary code of conduct written by astronomers for astronomers, and it survives on
goodwill.

So the honest state of global governance in one line: an unenforceable 1989 declaration, one
non-binding 1978 General Assembly decision, a scattering of national inboxes, and a treaty article
about contamination written for Moon rocks. \hi{We are not prepared. We have not even been told
that we are not prepared.}

\subsection{Why Nobody Wants the Job}

The absence is not an accident and it is not a cover-up; it is an incentive structure, and it is
worth spelling out because it predicts the future better than any leak.

A minister who funds a first-contact office buys ridicule today against a benefit that will
probably never arrive during their term. An international body that adopts a protocol has to decide
who speaks for Earth, which is a question no state wants settled while it is not the one speaking.
Any disclosure framework runs directly into national-security classification, because the sensors
that would detect the event are the same sensors whose capabilities are secret --- and that is the
real reason military UAP files stay closed, far more often than anything hiding in them. Meanwhile
the scientific community has spent seventy years protecting its funding from association with the
subject, so the constituency that would normally lobby for preparedness is the constituency with
the most to lose by asking.

Add those four pressures together and you get precisely the world we observe: strong detection,
zero preparation, and a public that would learn of a confirmed detection from a press conference
nobody had rehearsed.

\subsection{What Preparation Would Actually Look Like}

If the question is what governments should be \emph{doing}, then it is only fair to say what the
answer would look like, so that the reader can check it off as it does or does not happen. None of
the following requires believing anything at all about whether we are being visited.

A designated verification chain, so that a candidate detection is confirmed by named independent
facilities on a published timetable rather than by rumour. A single accountable international point
of announcement, seated in the UN system, because a confirmed detection announced simultaneously by
forty governments in forty framings is a disinformation event by construction. A standing
consultation rule on reply, giving legal effect to the one SETI principle that most needs it. A
planetary-protection and biocontainment regime extended from returning samples to any recovered
material of unknown origin, which Article IX of the Outer Space Treaty gestures at and no
implementing instrument covers. Pre-drafted public communication --- the same unglamorous work that
pandemic and asteroid planning already do, written in advance and rehearsed, because the first
twenty-four hours of an announcement decide whether the population responds or panics. And, plainly,
data release schedules so that the archive of anomalies is a scientific resource rather than a
classification problem.

That list is not exotic. It is the International Asteroid Warning Network with the nouns changed.
The institutional template already exists and has been agreed through COPUOS; it has simply never
been applied to this hazard. \hi{The barrier is not capability, and it is not law. It is that
nobody has asked the General Assembly for it since a small Caribbean state did in 1978.}

\subsection{And Then the Debate We Actually Have}

With the governance ledger on the table, the argument that scientists do conduct in public can be
seen for what it is: a debate about a decision that no institution is empowered to make.

\actualq[argument]{whether they exist.}{whether we should be transmitting.}

Passive SETI --- listening --- is uncontroversial. \textbf{METI}, Messaging to Extraterrestrial
Intelligence, deliberately broadcasting to selected targets, is not, and the dispute has produced one
of the few genuinely well-conducted debates in this whole field, largely because both sides are
arguing under total uncertainty and both know it.

The case against transmitting rests on asymmetry. Any species that receives our signal and is capable
of reaching us is, necessarily, far ahead of us; contact between parties of grossly unequal capability
has a poor historical record; and the decision is irreversible, unilateral, and taken on behalf of
everyone alive and everyone not yet born. This is Hawking's position, and David Brin's, and it is
essentially an argument about decision-making under uncertainty rather than about aliens: when the
downside is unbounded and the upside is speculative, and when no one has the standing to consent for
the species, the default is to wait.

The case for transmitting is also strong. We have been leaking radio, radar and television for a
century and cannot recall it, so the horse has left; military early-warning radar is by far our
loudest emission and is orders of magnitude more detectable than any deliberate message; a
civilisation able to cross interstellar distance could detect our atmosphere's oxygen and methane
spectroscopically without our help at all; and a galaxy in which everyone reasons like Brin is a
galaxy of permanent mutual silence, which may be a sufficient explanation for the Fermi paradox all
by itself. Alexander Zaitsev, who actually transmitted several messages from the Yevpatoria dish in
Ukraine, made this argument and then acted on it unilaterally --- which rather demonstrates the
procedural objection.

My own view, offered as a view: the asymmetry argument is correct in structure and probably
irrelevant in practice. We are already visible to anyone with better instruments than ours. The
METI debate is therefore less about safety than about consent, and on consent the objectors are
plainly right --- one radio astronomer with dish time should not be able to introduce the species.

And this is where the chapter's two founders part company, politely and in public. Drake, having sent
the Arecibo message himself, remained relaxed about transmission and sceptical that a receiving
civilisation would be either interested in or capable of finding us; Shostak has argued for a
deliberate, internationally agreed broadcast on the grounds that concealment is already impossible and
that the silence policy assumes a hostility for which there is no evidence. Hawking took the opposite
side on the same evidence. \hi{Three people who agreed completely on the arithmetic disagreed completely
on what to do about it}, which tells you the dispute was never really about astronomy --- it is about how
much weight an unquantifiable catastrophic risk should carry, and that is a question science does not
settle.

\begin{funfact}
The most reckless transmission ever sent was not a scientific message. In 2008 NASA broadcast the
Beatles song ``Across the Universe'' towards Polaris to mark an anniversary. In 2005 a snack food
company transmitted an advertisement towards a star system. Whatever the outcome of the METI debate,
the historical record will show that humanity's considered position was overruled by a marketing
department.
\end{funfact}

% ---------------------------------------------------------------
\section{Guerrilla Warfare}\label{sec:guerrilla}

Now the negative case, taken seriously rather than as entertainment. If contact went badly, what
would badly actually look like?

The first thing to discard is the cinema. Every invasion narrative from \emph{The War of the Worlds}
onward --- and Section~\ref{sec:warworlds} is about why that particular story has such a grip --- imagines
something like a human war: landings, resistance, a decisive engagement, a discovered weakness. This
is a category error. A conflict between civilisations separated by even a few thousand years of
technological development is not a war, in the way that a fight between an aircraft carrier and a
bronze-age fleet is not a war. Where the capability gap is total, there is no engagement to speak of.

Which brings us to the section title, because guerrilla warfare is the only form of resistance that
is even structurally available to the weaker party, and it is worth understanding why. Asymmetric
warfare works by refusing the contest the stronger side is optimised for: dispersing so there is no
target worth striking, blending into a population the aggressor cannot distinguish, imposing cost
over time rather than seeking a decision, and winning by making occupation more expensive than
withdrawal. Vietnam and Afghanistan are the standard illustrations, and both worked for the same
reason --- the stronger party wanted something that required holding ground and dealing with people.

And that is precisely the assumption that fails here. Guerrilla resistance depends on the adversary
needing the territory intact and the population present. Remove that need and the strategy has
nothing to grip. Worse, the classic motives for invasion mostly evaporate on inspection: water is
abundant in the outer solar system and easier to collect from a comet than from the bottom of a
gravity well defended by a biosphere; metals are more accessible in asteroids; and anything capable
of crossing interstellar space has energy sources that make ours irrelevant. Earth's only genuinely
rare feature is its biology --- and that is also the one thing that would make it hazardous, since
four billion years of local pathogens and a fully occupied biochemistry is a hostile environment,
not a resource. If there is a real risk, it is much more likely indifference than malice: the road
is not built through the anthill out of hatred.

\pullquote[Stephen Hawking]{``If aliens visit us, the outcome would be much as when Columbus landed in
America, which didn't turn out well for the Native Americans.''}

There is one final scenario in this direction that I think is underrated, and it is not military at
all. The most destructive contact event in human history did not involve weapons; it involved
microorganisms and information. Contact between human cultures has repeatedly destroyed the weaker
party through disease it had no immunity to and through the collapse of its own explanatory framework
--- the discovery that the world was not as it had been described. Both of those mechanisms are
available to a contact event that consists of nothing but a received message.

\figwrap[r]{0.40}{A U.S.\ psychological-operations soldier wiring a tape recorder to a loudspeaker in Vietnam. The hardware of the ``ghost tape'' campaign was this ordinary.}{https://commons.wikimedia.org/wiki/File:Psyops_soldier_sets_up_loudspeaker_during_2nd_Squadron,_11th_Cavalry_Regiment_MEDCAP_in_Xa_Hung_Loc.jpg}{fig:psyoploud}
And here I want to put a very concrete case on the table, because we have already run this
experiment on ourselves, in a small way, and the man who later hosted the world's largest
paranormal radio programme was in the country while it was happening. Art Bell --- whose show gets a
chapter of its own in Section~\ref{sec:artbell} --- served in the United States Air Force during the
Vietnam War as a medic. I have found nothing to suggest he had any involvement in what follows, and I
am not going to invent a connection; the point is simply that the war he served in produced the
cleanest illustration I know of the argument this chapter has been making about \hi{information as a
weapon}.

The operation was called \emph{Wandering Soul}, and its most famous artefact was a recording known to
the soldiers who played it as \emph{Ghost Tape Number Ten}. Vietnamese funeral custom holds that the
dead must be buried in their home ground; a man killed and left far from home becomes a spirit that
wanders, unquiet, unable to rest. American psychological-operations units --- the 6th Psychological
Operations Battalion among them, working with the U.S.\ Navy and with South Vietnamese speakers
brought in so that the accents would be right --- built a tape out of that belief. Buddhist funeral
music, echo and reverberation added in the studio, a child's voice calling ``Come home, Daddy'', and
the voice of a dead soldier addressing his former comrades:

\pullquote[\emph{Ghost Tape Number Ten}, translated]{``My body is gone. I am dead, my family. Tragic,
how tragic! My friends, I come back to let you know that I am dead\ldots\ I am in hell. Friends, while
you are still alive\ldots\ go home! Go home, my friends --- before it is too late!''}

They played it at night, from helicopters and from loudspeakers mounted on boats and trucks, into the
jungle, at men who were trying to sleep. \hi{No shot was fired by the tape.} The intended casualty was
morale, and the delivery mechanism was a piece of the enemy's own cosmology handed back to him after
dark.

What happened next is the part I find genuinely instructive, and it is not the part the story is
usually told for. It largely did not work. Soldiers who worked out that they were listening to a
loudspeaker did the obvious thing and shot at it --- which, as the Army's own assessment team noted
with some satisfaction, revealed their positions. The team that ran \emph{Wandering Soul} conceded
that the enemy ``realised what was going on'' and then declared the operation a success anyway,
without offering evidence for the claim. The Americans stopped using the tape in the early 1970s. A
variant intended for the Congo, in which recordings would impersonate angry local gods to keep
villagers at home, was prepared and never deployed.

So we have, on the record, a technologically superior party attempting to defeat a weaker one by
transmitting a supernatural message calibrated to that party's beliefs about death --- and mostly
failing, because the target population contained sceptics who noticed the loudspeaker. I do not think
the analogy to interstellar contact is exact, and I want to be careful with it: a civilisation
capable of crossing between stars would not be limited to our production values, and I am not
claiming that scepticism scales up as a defence. But the structure of the episode is exactly the
structure of the risk this chapter has been circling. \hi{The attack surface was not the body. It was
the cosmology.} If you want to know what a hostile transmission would actually try to do to a
species, \emph{Wandering Soul} is the closest thing to a controlled trial we have, and the trial was
run by us, on other human beings, using nothing but a tape recorder and a belief about the dead.

Which also means the lesson runs the other way, and I would rather say so plainly than end on a
frisson. The defence in Vietnam was not counter-technology. It was somebody in the dark deciding that
a voice claiming to be a ghost was more likely to be a speaker on a helicopter, and turning out to be
right. That is the same instrument this entire book has been trying to hand you.

% ---------------------------------------------------------------
\section{We Come in Peace}\label{sec:wecomepeace}

The optimistic case deserves the same seriousness, and it has better arguments than its detractors
admit.

Start with the selection effect. Any species that survives long enough to cross interstellar
distances has, by definition, solved the problems that would otherwise have killed it: it has
controlled the energies required for starflight without destroying itself, which is a filter humanity
has not yet passed. Carl Sagan's version of this argument is that longevity and aggression are
probably anti-correlated at the civilisational scale, because the technologies of interstellar travel
are the technologies of self-annihilation. A galaxy full of aggressive young civilisations is a galaxy
full of very short-lived ones. Whoever is still out there after a million years likely got that way
by being stable.

Second, motive. A civilisation with no material need of us has no reason for hostility, and the same
argument that undercut the invasion scenario in the last section cuts here too --- but in the
opposite direction it leaves curiosity standing. Information is the one thing Earth has that cannot
be obtained anywhere else: an independent four-billion-year biological experiment, a second data point
on how life arises and what it does. That is genuinely valuable, non-consumable, and best obtained by
leaving the specimen intact and talking to it.

Third, the practice already exists. There is a serious literature on how a benign contact might
actually proceed, and its most attractive proposal is that a sufficiently advanced civilisation would
not announce itself at all but leave material accessible on our terms --- what the field calls the
zoo hypothesis in its passive form, or in Iain Banks's phrasing, contact conducted at the pace of
the less capable party. The interesting feature of that model is that it is indistinguishable from
absence, which is intellectually inconvenient and worth admitting.

And fourth, the human institutional response might be better than we fear. When COSPAR wrote
planetary protection protocols, when the Antarctic Treaty set aside a continent for science, when the
Outer Space Treaty of 1967 declared that space is not subject to national appropriation, the species
managed --- under no pressure and with no enforcement --- to behave with a foresight that its general
reputation does not predict. Section~\ref{sec:rio}'s post-detection protocols are in that tradition.
They are unenforceable and they were still written down, by people trying to get it right in advance.
I find that more encouraging than most of what is in this book.

% ---------------------------------------------------------------
\section{The Problem With SETI}\label{sec:setiproblem}

And here is where I part company with the mainstream, because the standard summary of SETI's
sixty-year record is that we have looked and found nothing, and that summary is close to worthless.

The problem is radio, and specifically the fraction of the search space that radio SETI has actually
examined. Consider what a complete search would require: every direction on the sky, every frequency
from the low radio through the optical, every polarisation, every possible modulation scheme, and ---
this is the part that gets forgotten --- every \emph{moment in time}, because a signal you are not
pointed at when it arrives is a signal that did not happen as far as your data are concerned. Now
consider what has been done. A few thousand nearby stars observed attentively, mostly in a narrow
band around the hydrogen line, mostly for minutes at a time, at sensitivities that would generally
require a deliberate high-power beacon aimed at us to register.

Jill Tarter's own analogy, and she ran the search for decades, is that the totality of SETI's
coverage is comparable to dipping a drinking glass into the ocean, examining it, and concluding there
are no fish. The published estimates of the fraction of the plausible search space covered come out
at something like one part in $10^{18}$ --- a number so small that ``we have found nothing'' carries
almost no information about whether anything is there. Rounded to any practical precision, we have
searched zero.

Shostak's counter to his own colleague's analogy is worth setting out fairly, because it is the most
falsifiable prediction anyone in this field has made. His argument is that the coverage figure is a
snapshot of a process running on exponential hardware: receiver bandwidth, channel counts and
processing throughput have been doubling on the usual timescales, so the number of star systems SETI can
examine per decade grows accordingly. On that basis he has repeatedly said in print and in testimony to
Congress that if the optimistic end of Drake's equation is anywhere near right, a detection should
arrive within roughly two dozen years of speaking. I have watched him say a version of this for most of
my adult life, and the deadline keeps moving --- which he freely admits.

But notice what he did there. He attached a date to a claim, in public, knowing it could be checked
against him. \hi{The only person in this entire subject who has volunteered a falsifiable deadline is a
SETI astronomer, and it has been running against him for thirty years.} Compare that with the disclosure
industry's permanently imminent revelations, which never carry a date and are never scored when they
fail. Drake, characteristically, was more cautious: he revised his own equation's terms downward over his
lifetime as the data came in, and said in his last interviews that he expected the answer to arrive long
after he was gone. He was right about that much.

Then add the assumption problem on top, because coverage is not the only issue. Radio SETI assumes
that an advanced civilisation broadcasts, deliberately, at high power, continuously, omnidirectionally
or towards us specifically, using amplitude or frequency modulation on a narrow carrier, for a duration
overlapping our observation, from a distance where inverse-square losses still leave a detectable
signal. Every one of those is a guess derived from mid-twentieth-century human engineering. Our own
radio footprint has been shrinking for thirty years: we abandoned high-power analogue broadcast for
low-power digital, cable, fibre and tightly beamed satellite links, and a well-optimised digital
transmission is statistically indistinguishable from noise unless you know the encoding. A
civilisation like ours in 2026 would be substantially harder for SETI to detect than a civilisation
like ours in 1965. The window during which a technological species is loud in radio may be a century
out of a lifespan of millions of years --- and if so, radio SETI is looking for a signature that
barely exists.

\begin{funfact}
There is also the sheer arrogance of the frequency choice, which nobody likes to dwell on. We listen
at the hydrogen line because \emph{we} think it is the obvious universal channel. This is the
cosmic-scale equivalent of standing in a field shouting in English on the grounds that English is
widely spoken. The assumption may well be right. It is still an assumption, and it has consumed most
of the field's total observing time.
\end{funfact}

None of this is an argument that SETI is a waste. It is the most intellectually honest programme in
this entire book --- it states its hypothesis, it publishes null results, it does not inflate its
claims, and it has been doing rigorous science on almost no money for sixty years while the rest of
this subject was selling paperbacks. \hi{Two men are largely responsible for that culture}, and they are
the two this chapter opened with: Drake, who insisted the question be written as an equation that could
be argued about term by term, and Shostak, who spent three decades explaining the nulls to a public that
wanted announcements. My complaint is not about SETI. It is about how its results get
reported. Absence of evidence is evidence of absence only when the search was adequate to find the
thing, and Section~\ref{sec:allevidencebutnoproof}'s distinction between evidence and proof cuts both ways: a null result from an
instrument that covers a vanishing fraction of the possibility space is not a finding about the
universe. It is a finding about the instrument. Meanwhile the interesting work has moved --- optical
and infrared SETI looking for nanosecond laser pulses, technosignature searches for industrial
pollutants and waste heat in exoplanet atmospheres, and artefact searches within our own solar system,
where anything left behind would still be sitting there regardless of whether its makers are.

% ---------------------------------------------------------------
\section{Messages For The Future}\label{sec:messagesfuture}

Which brings the chapter, and in a sense the argument, to the Voyager Golden Record.

\figwrap[r]{0.42}{The Voyager Golden Record cover: pictographic playback instructions and a pulsar map giving Earth's position.}{https://commons.wikimedia.org/wiki/File:Voyager_Golden_Record_fx.png}{fig:goldenrecord}
In 1977 two spacecraft left Earth carrying identical gold-plated copper phonograph records (see Figure~\ref{fig:goldenrecord}),
compiled by a committee chaired by Carl Sagan, with Ann Druyan as creative director and Timothy Ferris
producing. Each carries 115 analogue-encoded images, greetings in 55 languages, a twelve-minute
sequence of the sounds of Earth --- surf, wind, thunder, birds, whales, a kiss, a mother and child, a
heartbeat, a Saturn V launch --- and ninety minutes of music: Bach, Beethoven, Mozart, Stravinsky,
Javanese gamelan, Senegalese percussion, Navajo night chant, Peruvian panpipes, Azerbaijani
balaban, Blind Willie Johnson's \emph{Dark Was the Night, Cold Was the Ground}, and Chuck Berry's
\emph{Johnny B. Goode}. The cover is engraved with instructions: how to play the record, the
rotational speed given in units of the hydrogen hyperfine transition, the pulsar map from the Pioneer
plaques, and a diagram of hydrogen itself as a key to the whole notation.

Both spacecraft are now in interstellar space. Voyager 1 crossed the heliopause in 2012, Voyager 2 in
2018. Their power sources will fail in the 2020s or early 2030s and after that they will drift,
silent, for a very long time --- the records were engineered to survive on the order of a billion
years, which is longer than the expected remaining habitability of the Earth's surface.

And here is the thing I want to leave you with. The probability that either record is ever found is,
realistically, indistinguishable from zero: they are moving slowly, they are not aimed at anything,
and interstellar space is almost entirely empty. Sagan knew that perfectly well. He said the record
would only be played if there were advanced spacefaring civilisations in interstellar space, and that
launching this bottle into the cosmic ocean said something very hopeful about life on this planet.

Which means the Golden Record was never really addressed outward. Consider what it contains: not our
weapons, our borders, our wars or our politics, but whale song, a mother greeting her child, a
Brandenburg concerto, and a Delta blues recording made by a blind street singer in 1927. That is a
species deciding what it would want to be judged by, and then committing the answer to gold and
putting it somewhere it can never be revised. The message that matters is not the one a stranger might
receive in a billion years. It is the one we sent ourselves in 1977, and the fact that when we
seriously imagined being seen from outside, this is what we chose to show.

Section~\ref{sec:mathlang} argued that mathematics gets you to ``we are both here and we both
count'' and no further. The Golden Record is the attempt to go further anyway --- to transmit not
information but \emph{what it is like to be us} --- knowing full well that nobody is listening. I
think that is the most human object ever built, and it is a fitting place to end a chapter dedicated
to a man who spent thirty years unable to speak and never once stopped trying to explain the universe
to anybody who would listen.

\begin{srcnotes}
\item Drake, F.\ (1961). Project Ozma. \emph{Physics Today}, 14(4).
\item Drake, F.\ \& Sobel, D.\ (1992). \emph{Is Anyone Out There? The Scientific Search for
Extraterrestrial Intelligence}.
\item Shostak, S.\ (2009). \emph{Confessions of an Alien Hunter: A Scientist's Search for
Extraterrestrial Intelligence}.
\item Shostak, S.\ (2014). Testimony before the U.S.\ House Committee on Science, Space and
Technology, hearing on Astrobiology and the Search for Life Beyond Earth, 21 May.
\item Overbye, D.\ (2022). Frank Drake, who led search for life on other planets, dies at 92.
\emph{The New York Times}, 3 September.
\item Almár, I.\ \& Tarter, J.\ (2011). The discovery of ETI as a high-consequence, low-probability
event. \emph{Acta Astronautica}, 68.
\item Loeb, A.\ (2025--2026). Essays and preprints on the anomalies of 3I/ATLAS; and
Breakthrough Listen radio observations of 3I/ATLAS (2026), reporting no technosignature and
ordinary cometary activity.
\item Forgan, D.\ et al.\ (2018). Rio 2.0: revising the Rio Scale for SETI detections.
\emph{International Journal of Astrobiology}.
\item Sagan, C., Drake, F.\ et al.\ (1978). \emph{Murmurs of Earth: The Voyager Interstellar
Record}.
\item Hawking, S.\ (2010). Remarks on extraterrestrial contact, \emph{Into the Universe with Stephen
Hawking}, Discovery Channel.
\item International Academy of Astronautics, Declaration of Principles Concerning Activities
Following the Detection of Extraterrestrial Intelligence (revised 2010).
\item ``Operation Wandering Soul'', Wikipedia --- campaign summary, the belief about unburied dead,
the construction of \emph{Ghost Tape Number Ten} and the Congo variant that was never used:
\srclink{https://en.wikipedia.org/wiki/Operation_Wandering_Soul}.
\item Friedman, H.\ A.\ (SGM, Ret.). The ``Wandering Soul'' tape of Vietnam. \emph{PsyWarrior} ---
first-hand account of the recording and its deployment:
\srclink{https://www.psywarrior.com/wanderingsoul.html}.
\item Watson, P.\ (1978). \emph{War on the Mind}. London: Hutchinson, pp.\ 410--411 --- on the
assessment team's claim of success and the prepared Congo operation.
\item Young, M.\ B.\ \& Buzzanco, R.\ (eds.) (2008). \emph{A Companion to the Vietnam War}.
Wiley, p.\ 442.
\item Humphrey, C.\ The ghostly legacies of America's war in Vietnam. \emph{Foreign Policy}.
\item Bell, A.\ --- U.S.\ Air Force service as a medic during the Vietnam War; see also the sources
listed for Chapter~\ref{ch:media}.
\end{srcnotes}
